
\section{Maximum principles} \label{maxprin}

In this appendix we list some maximum principles and their
consequences.  Our main source is \cite{Chow-Lu-Ni}, where
references to the original literature can be found.

The first type of maximum principle is a {\em weak} maximum
principle which says that under certain conditions, a spatial inequality
on the initial condition implies a time-dependent inequality at later times.

\begin{theorem}[maximum principles theorem]
  \label{thm:kl-maximum-principles-maximum-principles-theorem}
  \dcref{Theorem A.1}
  \group{Kleiner--Lott Analytic Appendices}
  \source{kleiner-lott}

Let $M$ be a closed manifold. Let $\{g(t)\}_{t \in [0,T]}$ be a smooth
one-parameter family of Riemannian metrics on $M$ and let
$\{X(t)\}_{t \in [0, T]}$ be
a smooth one-parameter family of vector fields on $M$. Let $F \: : \:
\R \times [0, T] \rightarrow \R$ be a Lipschitz function.  Suppose that
$u=u(x,t)$ is $C^2$-regular in $x$, $C^1$-regular in $t$ and
\begin{equation}
\frac{\partial u}{\partial t} \: \le \: \triangle_{g(t)} u \: + \: X(t) u \: + \: F(u, t).
\end{equation}
Let $\phi \: : \: [0,T] \rightarrow \R$ be the solution of
$\frac{d\phi}{dt} \: = \: F(\phi(t), t)$ with $\phi(0) = \alpha$. If
$u(\cdot, 0) \le \alpha$ then $u(\cdot, t) \le \phi(t)$ for all $t \in [0, T]$.
\end{theorem}

There are various noncompact versions of the weak maximum principle.
We state one here.

\begin{theorem}[estimates for curvature]
  \label{thm:kl-maximum-principles-estimates-curvature}
  \dcref{Theorem A.3}
  \group{Kleiner--Lott Analytic Appendices}
  \source{kleiner-lott}

Let $(M, g(\cdot))$ be a complete Ricci flow solution on the interval
$[0,T]$ with uniformly bounded curvature.  If $u=u(x,t)$ is a
$W^{1,2}_{loc}$ function that weakly satisfies
$\frac{\partial u}{\partial t} \: \le \: \triangle_{g(t)} u$, with
$u(\cdot, 0) \le 0$ and
\begin{equation}
\int_0^T \int_M e^{- \: c \: d_t^2(x, x_0)} \: u^2(x,s) \: dV(x) \: ds < \infty
\end{equation}
for some $c > 0$, then $u(\cdot,t) \le 0$ for all $t \in [0,T]$.
\end{theorem}

A {\em strong} maximum principle says that under certain conditions,
a strict inequality at a given time implies strict inequality at later
times and also slightly earlier times.  It does not require complete metrics.

\begin{theorem}[existence for maximum principles]
  \label{thm:kl-maximum-principles-existence-maximum-principles}
  \dcref{Theorem A.5}
  \group{Kleiner--Lott Analytic Appendices}
  \source{kleiner-lott}

Let $M$ be a connected manifold. Let $\{g(t)\}_{t \in [0,T]}$ be a smooth
one-parameter family of Riemannian metrics on $M$ and let
$\{X(t)\}_{t \in [0, T]}$ be
a smooth one-parameter family of vector fields on $M$. Let $F \: : \:
\R \times [0, T] \rightarrow \R$ be a Lipschitz function.  Suppose that
$u=u(x,t)$ is $C^2$-regular in $x$, $C^1$-regular in $t$ and
\begin{equation}
\frac{\partial u}{\partial t} \: \le \: \triangle_{g(t)} u \: + \: X(t) u \: + \: F(u, t).
\end{equation}
Let $\phi \: : \: [0,T] \rightarrow \R$ be a solution of
$\frac{d\phi}{dt} \: = \: F(\phi(t), t)$. If
$u(\cdot, t) \le \phi(t)$ for all $t \in [0, T]$ and
$u(x_0, t_0) < \phi(t)$ for some $x_0 \in M$ and $t_0 \in (0,T]$ then
there is some $\epsilon > 0$ so that $u(\cdot, t) < \phi(t)$ for
$t \in (t_0 - \epsilon, T]$.
\end{theorem}

A consequence of the strong maximum principle is a statement about
restricted holonomy for Ricci flow solutions with nonnegative curvature
operator $\Rm$.

\begin{theorem}[existence for curvature operators]
  \label{thm:kl-maximum-principles-existence-curvature-operators}
  \dcref{Theorem A.7}
  \group{Kleiner--Lott Analytic Appendices}
  \source{kleiner-lott}
 \label{splitting}
Let $M$ be a connected manifold. Let $\{g(t)\}_{t \in [0,T]}$ be a smooth
one-parameter family of Riemannian metrics on $M$ with nonnegative
curvature operator that satisfy the Ricci flow equation.  Then for each $t \in (0, T]$, the
image $\Image(\Rm_{g(t)})$ of the curvature operator
is a smooth subbundle of $\Lambda^2(T^* M)$ which is invariant under
spatial parallel translation. There is a sequence of times
$0 = t_0 < t_1 < \ldots < t_k = T$ such that for each
$1 \le i \le k$,
$\Image(\Rm_{g(t)})$ is a Lie subalgebra of $\Lambda^2(T^*_m M)
\cong o(n)$ that is independent of $t$ for $t \in (t_{i-1}, t_i]$.
Furthermore, $\Image(\Rm_{g(t_i)}) \subset \Image(\Rm_{g(t_{i+1})})$.
\end{theorem}

In particular, under the hypotheses of Theorem \ref{splitting}, a
local isometric splitting at a given time implies a local
isometric splitting at earlier times.


\section{$\phi$-almost nonnegative curvature}
\label{phiappendix}

In three dimensions,
the Ricci flow equation implies that
\begin{equation} \label{3dR}
\frac{dR}{dt} \: = \:
\triangle R \: + \: \frac{2}{3} R^2 \: +
2 \: |R_{ij} - \frac{R}{3} g_{ij}|^2 .
\end{equation}
The maximum principle of Appendix \ref{maxprin} implies that
if $(M, g(\cdot))$ is a Ricci flow solution defined for
$t \in [0, T)$, with complete time slices and bounded curvature
on compact time intervals, then
\begin{equation} \label{Rlowerbound}
(\inf R)(t) \: \ge \:
\frac{(\inf R)(0)}
{1 - \frac23 t (\inf R)(0)}.
\end{equation}
In particular,
$tR(\cdot, t) \: > \: - \: \frac{3}{2}$ for all $t \ge 0$
(compare \cite[Section 2]{Hamiltonn}).

Recall that the curvature operator is an operator on $2$-forms.
We follow the usual Ricci flow convention that if a manifold
has constant sectional curvature $k$ then its curvature operator
is multiplication by $2k$. In general, the trace of the
curvature operator equals the scalar curvature.

In three dimensions, having nonnegative
curvature operator is equivalent to having nonnegative
sectional curvature.  Each eigenvalue of the curvature operator
is twice a sectional curvature.

Hamilton-Ivey pinching, as given
in \cite[Theorem 4.1]{Hamiltonn}, says the following.

Assume that at $t = 0$ the eigenvalues $\lambda_1 \le \lambda_2 \le
\lambda_3$ of the
curvature operator at each point satisfy $\lambda_1 \ge -1$.
(One can
always achieve this by rescaling. Note that it implies
$R(\cdot, t) \ge - \frac32 \frac{1}{t+\frac12}$.)
Given a point $(x, t)$, put $X \: = \: - \lambda_1$.
If $X > 0$ then
\begin{equation}
R(x, t) \: \ge \: X \left( \ln X + \ln(1+t) - 3 \right),
\end{equation}
or equivalently,
\begin{equation} \label{estimate}
tR(x, t) \: \ge \: tX \left( \ln(tX) + \ln(\frac{1+t}{t}) - 3 \right).
\end{equation}

\begin{definition}[positivity for curvature operators]
  \label{def:kl-almost-nonnegative-curvature-positivity-curvature-operators}
  \dcref{Definition B.5}
  \group{Kleiner--Lott Analytic Appendices}
  \source{kleiner-lott}
 \label{HamIvey}
Given $t\geq 0$, a Riemannian $3$-manifold $(M,g)$ satisfies the
{\em  time-$t$ Hamilton-Ivey pinching condition} if for every $x\in M$,
if $\la_1\leq \la_2\leq \la_3$ are the eigenvalues of the curvature operator at $x$,
then either

$\bullet$ $\la_1\geq 0$, i.e. the curvature is nonnegative at $x$,

or

$\bullet$ If $\la_1 < 0$ and $X \: = \: - \lambda_1$ then
$tR(x)\geq tX \left(\ln (tX) + \ln \left( \frac{1+t}{t} \right)-3
\right)$.
\end{definition}
This condition has the following monotonicity property:

\begin{lemma}[monotonicity for scalar curvature]
  \label{lem:kl-almost-nonnegative-curvature-monotonicity-scalar-curvature}
  \dcref{Lemma B.6}
  \group{Kleiner--Lott Analytic Appendices}
  \source{kleiner-lott}

\label{lempinchingmonotonic}
Suppose that $\Rm$ and $\Rm'$
are $3$-dimensional  curvature operators whose scalar curvatures
and first eigenvalues satisfy
$R'\geq R\geq 0$ and $\la_1'\geq \la_1$.
If $\Rm$ satisfies the time-$t$ Hamilton-Ivey pinching condition then so does
$\Rm'$.
\end{lemma}
\begin{proof}
We may assume that  $\la_1'<0$ and
$\log (tX') \: + \: \ln \left( \frac{1+t}{t} \right) \: - \: 3 \: > \: 0$,
since otherwise the condition will be satisfied
(because $R'>0$ by hypothesis).  The function
\begin{equation}
Y\mapsto tY \left( \ln(t Y) \: + \: \ln \left( \frac{1+t}{t} \right) \: - \: 3 \right)
\end{equation}
is monotone increasing on the interval on which
$\ln(t Y) \: + \: \ln \left( \frac{1+t}{t} \right) \: - \: 3$
is nonnegative, so
$tX \left( \ln(t X) \: + \: \ln \left( \frac{1+t}{t} \right) \: - \: 3 \right)
\: \ge \:
tX^\prime \left( \ln(t X^\prime) \: + \: \ln \left( \frac{1+t}{t} \right) \: - \: 3 \right)$.
Hence $\Rm'$ satisfies the pinching condition too.
\end{proof}


The content of the pinching equation is that for any $s \in \R$,
if $t R(\cdot, t) \: \le \:
s$ then there is a lower bound $t \Rm (\cdot, t) \: \ge \: \const
(s, t)$. Of course, this is a vacuous statement
if $s \: \le \: - \: \frac{3}{2}$.

Using equation (\ref{estimate}), we can find a positive
function $\Phi \in C^\infty(\R)$ such that \\
1. $\Phi$ is nondecreasing. \\
2. For $s > 0$, $\frac{\Phi(s)}{s}$ is decreasing. \\
3. For large $s$,
$\Phi(s) \sim \frac{s}{\ln s}$. \\
4. For all $t$,
\begin{equation} \label{phi}
\Rm(\cdot, t) \: \ge \: - \: \Phi(R(\cdot, t)).
\end{equation}

This bound has the most consequence when $s$ is large.

We note that for the original unscaled Ricci flow
solution, the precise
bound that we obtain
depends on $t_0$ and
the time-zero metric, through its lower curvature bound.

\section{Ricci solitons} \label{Riccisolitons}

Let $\{V(t)\}$ be a time-dependent family of vector fields on a manifold $M$.
The solution to the equation
\begin{equation}
\frac{dg}{dt} \: = \: {\mathcal L}_{V(t)} g
\end{equation}
is
\begin{equation}
g(t) \: = \: \phi^{-1}(t)^* g(t_0)
\end{equation}
where $\{ \phi(t) \}$ is the $1$-parameter group of diffeomorphisms generated
by $- V$, normalized by $\phi(t_0) \: = \: \Id$.
(If $M$ is noncompact then we assume that $V$ can be integrated.
The reason for the funny signs is that if a $1$-parameter family of
diffeomorphisms $\eta(t)$ is generated by vector fields $W(t)$ then
${\mathcal L}_{W(t)} \: = \: \eta^{-1}(t)^* \:
\frac{d\eta(t+\epsilon)^*}{d\epsilon} \Big|_{\epsilon = 0} \: = \: - \:
\frac{d\eta^{-1}(t+\epsilon)^*}{d\epsilon} \Big|_{\epsilon = 0} \: \eta(t)^*$.)

The equation for a {\em steady soliton} is
\begin{equation}
2 \Ric \: + \: {\mathcal L}_V g \: = \: 0,
\end{equation}
where $V$ is a time-independent
vector field.
The corresponding Ricci flow is given by
\begin{equation}
g(t) \: = \: \phi^{-1}(t)^* g(t_0),
\end{equation}
where $\{ \phi(t) \}$ is the $1$-parameter group of diffeomorphisms generated
by $- V$. (Of course, in this case
$\{ \phi^{-1}(t) \}$ is the $1$-parameter group of diffeomorphisms generated
by $V$.)

A {\em gradient steady soliton}
satisfies the equations
\begin{align} \label{steady}
\frac{\partial g_{ij}}{\partial t} \: & = \: - \: 2 R_{ij} \: = \:
2 \nabla_i \nabla_j f, \\
\frac{\partial f}{\partial t} \: & = \: |\nabla f|^2. \notag
\end{align}
It follows from (\ref{steady}) that
\begin{equation}
\frac{\partial}{\partial t} \left( g^{ij} \: \partial_j f \right) \: = \:
2 \: R^{ij} \: \nabla_j f \: + \: g^{ij} \nabla_j |\nabla f|^2 \: = \:
- 2 \: (\nabla^i \nabla^j f) \: \nabla_j f \: + \: \nabla^i
|\nabla f|^2 \: = \: 0,
\end{equation}
showing that $V \: = \: \nabla f$ is indeed constant in $t$.
The solution to (\ref{steady}) is
\begin{align} \label{steadysoln}
g(t) \: & = \: \phi^{-1}(t)^* g(t_0), \\
f(t) \: & = \: \phi^{-1}(t)^* f(t_0). \notag
\end{align}

Conversely, given a metric $\widehat{g}$ and a function $\widehat{f}$
satisfying
\begin{equation}
\widehat{R}_{ij} \: + \: \widehat{\nabla}_i \widehat{\nabla}_j \widehat{f}
\: = \: 0,
\end{equation}
put $V \: = \: \widehat{\nabla} \widehat{f}$.
If we define $g(t)$ and $f(t)$ by
\begin{align}
g(t) \: & = \: \phi^{-1}(t)^* \widehat{g}, \\
f(t) \: & = \: \phi^{-1}(t)^* \widehat{f} \notag
\end{align}
then they satisfy (\ref{steady}).


A solution to (\ref{steady}) satisfies
\begin{equation}
\frac{\partial f}{\partial t} \: = \: |\nabla f|^2 \: - \: \triangle f
\: - \: R,
\end{equation}
or
\begin{equation}
\frac{\partial}{\partial t} \: e^{-f} \: = \: - \: \triangle e^{-f}
\: + \: R \: e^{-f}.
\end{equation}
This perhaps motivates Perelman's use of the backward heat
equation (\ref{backwards}).


A {\em shrinking soliton} lives on a time interval $(-\infty, T)$.
For convenience, we take $T \: = \: 0$. Then the equation is
\begin{equation}
2 \Ric \: + \: {\mathcal L}_V g \: + \: \frac{g}{t} \: = \: 0.
\end{equation}
The vector field $V = V(t)$ satisfies
$V(t) \: = \: - \: \frac{1}{t} \: V(-1)$.
The corresponding Ricci flow is given by
\begin{equation} \label{shrinkingflow}
g(t) \: = \: - \: t \: \phi^{-1}(t)^* g(-1),
\end{equation}
where $\{ \phi(t) \}$ is the $1$-parameter group of diffeomorphisms generated
by $- V$, normalized by $\phi(-1) \: = \: \Id$.


A {\em gradient shrinking soliton} satisfies the equations
\begin{align} \label{shrinking}
\frac{\partial g_{ij}}{\partial t} \: & = \: - \: 2 R_{ij} \: = \:
2 \nabla_i \nabla_j f \: + \: \frac{g_{ij}}{t}, \\
\frac{\partial f}{\partial t} \: & = \: |\nabla f|^2. \notag
\end{align}
It follows from (\ref{shrinking}) that
$V \: = \: \nabla f$ satisfies $V(t) \: = \: - \: \frac{1}{t} \: V(-1)$.
The solution to (\ref{shrinking}) is
\begin{align} \label{shrinkingsoln}
g(t) \: & = \: - \: t \: \phi^{-1}(t)^* g(-1), \\
f(t) \: & = \: \phi^{-1}(t)^* f(-1). \notag
\end{align}

Conversely, given a metric $\widehat{g}$ and a function $\widehat{f}$
satisfying
\begin{equation}
\widehat{R}_{ij} \: + \: \widehat{\nabla}_i \widehat{\nabla}_j \widehat{f}
\: - \: \frac{1}{2} \:
\widehat{g}
\: = \: 0,
\end{equation}
put $V(t) \: = \: - \: \frac{1}{t} \: \widehat{\nabla} \widehat{f}$.
If we define $g(t)$ and $f(t)$ by
\begin{align}
g(t) \: & = \: - \: t \: \phi^{-1}(t)^* \widehat{g}, \\
f(t) \: & = \: \phi^{-1}(t)^* \widehat{f} \notag
\end{align}
then they satisfy (\ref{shrinking}).

An {\em expanding soliton} lives on a time interval $(T, \infty)$.
For convenience, we take $T \: = \: 0$. Then the equation is
\begin{equation}
2\Ric \: + \: {\mathcal L}_V g \: + \: \frac{g}{t} \: = \: 0.
\end{equation}
The vector field $V = V(t)$ satisfies
$V(t) \: = \: \frac{1}{t} \: V(1)$.
The corresponding Ricci flow is given by
\begin{equation}
g(t) \: = \: t \: \phi^{-1}(t)^* g(1),
\end{equation}
where $\{ \phi(t) \}$ is the $1$-parameter group of diffeomorphisms generated
by $- V$, normalized by $\phi(1) \: = \: \Id$.

A {\em gradient expanding soliton} satisfies the equations
\begin{align} \label{expanding}
\frac{\partial g_{ij}}{\partial t} \: & = \: - \: 2 R_{ij} \: = \:
2 \nabla_i \nabla_j f \: + \: \frac{g_{ij}}{t}, \\
\frac{\partial f}{\partial t} \: & = \: |\nabla f|^2. \notag
\end{align}
It follows from (\ref{expanding}) that
$V \: = \: \nabla f$ satisfies $V(t) \: = \: \frac{1}{t} \: V(1)$.
The solution to (\ref{expanding}) is
\begin{align} \label{expandingsoln}
g(t) \: & = \: t \: \phi^{-1}(t)^* g(1), \\
f(t) \: & = \: \phi^{-1}(t)^* f(1). \notag
\end{align}

Conversely, given a metric $\widehat{g}$ and a function $\widehat{f}$
satisfying
\begin{equation}
\widehat{R}_{ij} \: + \: \widehat{\nabla}_i \widehat{\nabla}_j \widehat{f}
\: + \: \frac{1}{2} \:
\widehat{g}
\: = \: 0,
\end{equation}
put $V(t) \: = \: \frac{1}{t} \: \widehat{\nabla} \widehat{f}$.
If we define $g(t)$ and $f(t)$ by
\begin{align}
g(t) \: & = \: t \: \phi^{-1}(t)^* \widehat{g}, \\
f(t) \: & = \: \phi^{-1}(t)^* \widehat{f} \notag
\end{align}
then they satisfy (\ref{expanding}).

Obvious examples of solitons are given by Einstein metrics, with
$V = 0$. Any steady or expanding soliton on a closed
manifold comes from an Einstein metric.  Other examples of solitons
(see \cite[Chapter 2]{Chow-Knopf})
are : \\ \\
1. (Gradient steady soliton) The cigar soliton on $\R^2$ and the Bryant soliton
on $\R^3$. \\
2. (Gradient shrinking soliton) Flat $\R^n$ with $f \: = \: - \:
\frac{|x|^2}{4t}$. \\
3. (Gradient shrinking soliton) The shrinking cylinder
$\R \times S^{n-1}$ with
$f \: = \: - \: \frac{x^2}{4t}$, where $x$ is the coordinate on $\R$. \\
4. (Gradient shrinking soliton) The Koiso soliton on $\C P^2 \#
\overline{\C P^2}$.

\section{Local derivative estimates} \label{applocalder}

\begin{theorem}[existence for local derivative estimates]
  \label{thm:kl-local-derivative-estimates-existence-local-derivative-estimates}
  \dcref{Theorem D.1}
  \group{Kleiner--Lott Analytic Appendices}
  \source{kleiner-lott}
 \label{localder}
For any $\alpha, K, K^\prime, l \: \ge \: 0$ and $m, n \in \Z^+$,
there is some $C \: = \: C(\alpha, K, K^\prime, l,m, n)$ with the
following property.
Given $r > 0$,
suppose that $g(t)$ is a Ricci flow solution
for $t \in [0, \overline{t}]$, where $0 \: < \:
\overline{t} \: \le \:
\frac{\alpha r^2}{K}$, defined on an open neighborhood $U$ of a point
$p \in M^n$. Suppose that $\overline{B(p, r, 0)}$ is a compact
subset of $U$, that
\begin{equation}
|\Rm(x, t)| \: \le \: \frac{K}{r^{2}}
\end{equation}
for all $x \in U$ and $t \in [0, \overline{t}]$, and that
\begin{equation}
|\nabla^\beta \Rm(x, 0)| \: \le \: \frac{K^\prime}{r^{|\beta|+2}}
\end{equation}
for all $x \in U$ and $|\beta| \: \le \: l$.
Then
\begin{equation}
|\nabla^\beta \Rm(x, t)| \: \le \:
\frac{C}{
r^{|\beta|+2} \: \left( \frac{t}{r^2} \right)^{\frac{\max(m-l,0)}{2}}
}
\end{equation}
for all $x \in B \left(p, \frac{r}{2}, 0 \right)$,
$t \in (0, \overline{t}]$ and $|\beta| \: \le \: m$.

In particular,
$|\nabla^\beta \Rm(x, t)| \: \le \:
\frac{C}{r^{|\beta|+2}}$ whenever
$|\beta| \: \le \: l$.
\end{theorem}

The main case $l = 0$ of Theorem \ref{localder} is due to Shi
\cite{Shi2}. The extension to $l \: \ge \: 0$ appears in
\cite[Appendix B]{Lu-Tian}.

\section{Convergent subsequences of Ricci flow solutions}
\label{subsequence}

\begin{theorem}[existence for convergent subsequences of ricci flow solutions]
  \label{thm:kl-convergent-subsequences-ricci-flow-solutions-existence-convergent-subsequences-ricci-flow-solutions}
  \dcref{Theorem E.1}
  \group{Kleiner--Lott Analytic Appendices}
  \source{kleiner-lott}

\label{accumulate}
Given $r_0 \in (0, \infty]$,
let $\{g_i(t)\}_{i=1}^\infty$ be a sequence of Ricci flow solutions
on connected pointed manifolds $(M_i, m_i)$,
defined for $t \in (A, B)$ with $-\infty \le A < 0 < B \le \infty$.
We assume that for all $i$, $M_i$ equals the time-zero ball $B_0(m_i, r_0)$
and for all $r \in (0, r_0)$, $\overline{B_0(m_i, r)}$ is compact.
Suppose that the following
two conditions are satisfied : \\
1. For each $r \in (0, r_0)$ and each compact interval $I \subset (A,B)$,
there is an $N_{r,I} < \infty$ so that for all $t \in I$ and all $i$,
$\sup_{B_0(m_i, r) \times I} |\Rm(g_i)| \le N_{r,I}$, and \\
2. The
time-$0$ injectivity radii
$\{\inj(g_i(0))(m_i)\}_{i=1}^\infty$
are uniformly bounded below by a positive number.

Then after passing to a subsequence, the solutions converge smoothly
to a
Ricci flow solution $g_\infty(t)$ on a connected pointed manifold
$(M_\infty, m_\infty)$, defined for
$t \in (A, B)$, for which $M_\infty = B_0(m_\infty, r_0)$ and
$\overline{B_0(m_\infty, r)}$ is compact for all $r \in (0, r_0)$.
That is, for any compact interval $I \subset (A, B)$
and any $r < r_0$,
there are pointed
time-independent diffeomorphisms $\phi_{r,i} \: : \:
B_0(m_\infty, r) \rightarrow B_0(m_i,r)$ so that
$\{(\phi_{r,i} \times \Id)^* g_i\}_{i=1}^\infty$  converges smoothly
to $g_\infty$ on $B_0(m_\infty, r) \times I$.
\end{theorem}

Given the sectional curvature bounds, the lower
bound on the injectivity radii is equivalent to a lower bound
on the volumes of balls around $m_i$
\cite[Theorem 4.7]{Cheeger-Gromov-Taylor}.
Theorem \ref{accumulate} is a slight generalization of
\cite[Main Theorem]{Hamiltonnn}, in which $r_0 = \infty$ and
$N_{r,I}$ is independent of $r$; see Corollary \ref{accumulate3}
below.

There are many variants of the theorem with alternative hypotheses.
One can replace the interval $(A,B)$ with an interval $(A, B]$,
$-\infty \le A < 0 \le B < \infty$.
One can also replace the interval $(A, B)$ with an interval
$[A, B)$, $-\infty < A< 0 < B
\le \infty$, if in addition one has uniform time-$A$ bounds
$\sup_{B_A(m_i, r)} |\nabla^j \Rm(g_i(A))| \: \le \: C_{r, j}$.
Then using Appendix \ref{applocalder},
one gets smooth convergence to a limit solution $g_\infty$
on the time interval $[A, B)$.
(Without the time-$A$ bounds one would only get
$C^0$-convergence on $[A, B)$ and $C^\infty$-convergence
on $(A, B)$.)
There is a similar statement if
one replaces the interval $(A, B)$ with an interval
$[A, B]$, $-\infty < A< 0 \le B
< \infty$. There is a version in which balls are replaced
by annuli. One can generalize the hypotheses to allow for
an $r$-dependent time interval.

In the setting of Theorem \ref{accumulate}, suppose that
$r_0 = \infty$. Then the time-zero slice $(M_\infty, m_\infty, g_\infty(0))$
is complete, but it does not immediately follow that the
other time slices are complete.  We now give a
condition which will guarantee completeness, and which will be
sufficient for our purposes.

\begin{corollary}[existence for convergent subsequences of ricci flow solutions]
  \label{cor:kl-convergent-subsequences-ricci-flow-solutions-existence-convergent-subsequences-ricci-flow-solutions}
  \dcref{Corollary E.2}
  \group{Kleiner--Lott Analytic Appendices}
  \source{kleiner-lott}

\label{accumulate2}
Let $\{g_i(t)\}_{i=1}^\infty$ be a sequence of Ricci flow solutions
on connected pointed manifolds $(M_i, m_i)$,
defined for $t \in (A, 0]$ with $-\infty \le A < 0$.
Suppose that each time-zero slice $(M_i, m_i, g_i(0))$ is complete.
Suppose that the following
three conditions are satisfied : \\
1. For each $r \in (0, \infty)$ and each compact interval $I \subset (A,0]$,
there is an $N_{r,I} < \infty$ so that for all $i$,
$\sup_{B_0(m_i, r) \times I} |\Rm(g_i)| \le N_{r,I}$,  \\
2. For each compact interval $I \subset (A, 0]$, there  is some
$N^\prime_I \in (0, \infty)$ with the following property : for each
$r \in (0, \infty)$, there is some $J_{r,I} \in \Z^+$ so that whenever
$i \ge J_{r,I}$, we have
$\Ric(g_i) \ge - N^\prime_I g_i$ on $B_0(m_i, r) \times I$, and \\
3. The
time-$0$ injectivity radii
$\{\inj(g_i(0))(m_i)\}_{i=1}^\infty$
are uniformly bounded below by a positive number.

Then after passing to a subsequence, the solutions converge smoothly
to a
Ricci flow solution $g_\infty(t)$ on a connected pointed manifold
$(M_\infty, m_\infty)$, defined for
$t \in (A, 0]$, with complete time slices.
\end{corollary}
\begin{proof}
  \uses{rem:kl-length-distortion-estimates-structural-remark-length-distortion-estimates,thm:kl-convergent-subsequences-ricci-flow-solutions-existence-convergent-subsequences-ricci-flow-solutions}
Let $(M_\infty, m_\infty, g_\infty(\cdot))$ be constructed as in Theorem
\ref{accumulate}. Then on each compact time interval $I \subset (A, 0]$,
we have
$\Ric(g_\infty) \ge - N^\prime_I g_\infty$. By (\ref{similar}),
if $t \in I$ then for any $m_0, m_1 \in M_\infty$, we have
\begin{equation} \label{similar2}
\dist_t(m_0, m_1) \ge e^{- N^\prime_I t} \dist_0(m_0, m_1).
\end{equation}

Suppose that $\{m_j\}_{j=1}^\infty$ is a Cauchy sequence
in $(M_\infty, g_\infty(t))$. From (\ref{similar2}), it is also
a Cauchy sequence in $(M_\infty, g_\infty(0))$, and so has a
subsequence, which we relabel as $\{m_j\}_{j=1}^\infty$, that
converges to some $m^\prime$ in $(M_\infty, g_\infty(0))$.
However, for small $\epsilon > 0$, the restriction of the
identity map $(M_\infty, g_\infty(0)) \rightarrow
(M_\infty, g_\infty(t))$ to
$B_0(m^\prime, \epsilon)$ is biLipschitz.  It follows that
$\lim_{j \rightarrow \infty} m_j = m^\prime$ in
$(M_\infty, g_\infty(t))$.
\end{proof}

There is an obvious analog to Corollary \ref{accumulate2} in
which we assume that the Ricci flows are defined for $[0, B)$ and
for any compact interval $I \subset [0, B)$,
there is an upper bound $\Ric(g_i) \le N^\prime_I g_i$
on $B_0(m_i, r) \times I$ for large $i$.
If we assume double-sided curvature bounds
then the statement is as follows.

\begin{corollary}[existence for convergent subsequences of ricci flow solutions]
  \label{cor:kl-convergent-subsequences-ricci-flow-solutions-existence-convergent-subsequences-ricci-flow-solutions-3}
  \dcref{Corollary E.4}
  \group{Kleiner--Lott Analytic Appendices}
  \source{kleiner-lott}

\label{accumulate3}
Let $\{g_i(t)\}_{i=1}^\infty$ be a sequence of Ricci flow solutions
on connected pointed manifolds $(M_i, m_i)$,
defined for $t \in (A, B)$ with $-\infty \le A < 0 < B \le \infty$.
We assume that for all $i$, the time slice
$(M_i, m_i, g_i(0))$ is complete.
Suppose that the following
two conditions are satisfied : \\
1. For each compact interval $I \subset (A,B)$,
there is an $N_{I} < \infty$ with the following property :
for each $r \in (0, \infty)$, there is some $J_{r,I} \in \Z^+$
so that whenever $i \ge J_{r,I}$, we have
$\sup_{B_0(m_i, r) \times I} |\Rm(g_i)| \le N_{I}$, and \\
2. The
time-$0$ injectivity radii
$\{\inj(g_i(0))(m_i)\}_{i=1}^\infty$
are uniformly bounded below by a positive number.

Then after passing to a subsequence, the solutions converge smoothly
to a
Ricci flow solution $g_\infty(t)$ on a connected pointed manifold
$(M_\infty, m_\infty)$, defined for
$t \in (A, B)$, with complete time slices.
\end{corollary}

In the case when $J_{r,I} = 1$ for all $r$ and $I$,
i.e. $\sup_{M_i \times I} |\Rm(g_i)| \le N_{I}$,
Corollary \ref{accumulate3} is the same as
\cite[Main Theorem]{Hamiltonnn}.

\section{Harnack inequalities for Ricci flow} \label{appharnack}

We first recall the statement of the matrix Harnack inequality.
Put
\begin{align}
P_{abc} \: & = \: \nabla_a R_{bc} \: - \: \nabla_b R_{ac}, \\
M_{ab} \: & = \: \triangle R_{ab} \: - \: \frac12 \: \nabla_a \nabla_b
R \: + \: 2 \: R_{acbd} R_{cd} \: - \: R_{ac} R_{bc} \: + \:
\frac{R_{ab}}{2t}. \notag
\end{align}
Given a $2$-form $U$ and a $1$-form $W$, put
\begin{equation}
Z(U, W) \: = \: M_{ab} W_a W_b \: + \: 2 \: P_{abc} U_{ab} W_c
\: + \: R_{abcd} U_{ab} U_{cd}.
\end{equation}
Suppose that we have a Ricci flow for $t > 0$
on a complete manifold with bounded
curvature on each
compact time interval
and nonnegative curvature operator.
Hamilton's matrix Harnack inequality says that for all $t > 0$ and all
$U$ and $W$, $Z(U, W) \: \ge \: 0$ \cite[Theorem 14.1]{Hamilton}.

Taking
$W_a \: = \: Y_a$ and $U_{ab} \: = \: (X_a Y_b \: - \:
Y_a X_b)/2$ and
using the fact that
\begin{equation} \label{Rict}
\Ric_t(Y, Y) \: = \: (\triangle R_{ab}) Y^a Y^b \: + \:
2 \: R_{acbd} R_{cd} Y^a Y^b \: - \: 2 \: R_{ac} R_{bc} Y^a Y^b,
\end{equation}
we can write $2Z(U, W) \: = \: H(X, Y)$, where
\begin{align} \label{Hexpression}
H(X, Y) \: = \: & - \: \Hess_R (Y,Y) \: - \: 2 \langle R(Y, X)Y, X \rangle
\: +\: 4 \left( \nabla_X \Ric(Y, Y) \: - \: \nabla_Y \Ric(Y, X) \right) \\
&  + \: 2 \Ric_t(Y, Y) \: + \: 2 \big| \Ric(Y, \cdot) \big|^2
\: + \: \frac{1}{t} \: \Ric(Y, Y). \notag
\end{align}
Substituting the elements of an orthonormal basis $\{e_i\}_{i=1}^n$ for
$Y$ and summing over $i$ gives
\begin{align}
\sum_i H(X, e_i) \: = & \: - \: \triangle R \: + \: 2 \Ric(X, X) \: +  \\
& 4 (\langle \nabla R, X \rangle
 \: - \: \sum_i \nabla_{e_i} \Ric(e_i, X)) \: + \:
2 \sum_i \Ric_t(e_i, e_i) \: + \: 2 |\Ric|^2 \: + \: \frac{1}{t} \: R.
\notag
\end{align}
Tracing the second Bianchi identity gives
\begin{equation}
\sum_i \nabla_{e_i} \Ric(e_i, X) \: = \:
\frac{1}{2} \: \langle \nabla R, X \rangle.
\end{equation}
From (\ref{Rict}),
\begin{equation}
\sum_i \Ric_t(e_i, e_i) \: = \: \triangle R.
\end{equation}
Putting this together gives
\begin{equation} \label{together}
\sum_i H(X, e_i) \: = \: H(X),
\end{equation}
where
\begin{equation} \label{Heqn}
H(X) \: = \: R_t \: + \: \frac{1}{t} \: R \: + \:
2 \langle \nabla R, X \rangle \:
+ \: 2 \Ric(X, X).
\end{equation}
We obtain Hamilton's trace Harnack
inequality, saying that $H(X) \: \ge \: 0$ for all $X$.

In the rest of this section we assume that
the solution is defined for all $t \in (- \infty, 0)$.
Changing the origin point of time, we have
\begin{equation}
R_t \: + \: \frac{1}{t- t_0} \: R \: + \:
2 \langle \nabla R, X \rangle \:
+ \: 2 \Ric(X, X) \: \ge \: 0
\end{equation}
whenever $t_0 \: \le \: t$. Taking $t_0 \rightarrow - \infty$ gives
\begin{equation}
R_t \: + \:
2 \langle \nabla R, X \rangle \:
+ \: 2 \Ric(X, X) \: \ge \: 0
\end{equation}
In particular, taking $X = 0$ shows that the scalar curvature is
nondecreasing in $t$ for any ancient solution with nonnegative
curvature operator, assuming again that the metric is complete
on each time slice with bounded curvature
on each compact time interval.
More generally,
\begin{equation}
0 \: \le \: R_t \: + \:
2 \langle \nabla R, X \rangle \:
+ \: 2 \Ric(X, X) \: \le \: R_t \: + \:
2 \langle \nabla R, X \rangle \:
+ \: 2 R \langle X, X \rangle.
\end{equation}
If $\gamma \: : \: [t_1, t_2]
\rightarrow M$ is a curve parametrized by $s$ then taking
$X \: = \: \frac12 \: \frac{d\gamma}{ds}$ gives
\begin{equation}
\frac{d R(\gamma(s), s)}{ds}  \: = \:
R_t(\gamma(s), s) \: + \: \left\langle \frac{d\gamma}{ds}, \nabla R
\right\rangle
\: \ge \: - \: \frac12 \: R \: \left\langle \frac{d\gamma}{ds}, \frac{d\gamma}{ds}
\right\rangle.
\end{equation}
Integrating $\frac{d \ln R(\gamma(s),s)}{ds}$ with respect to $s$ and using the fact that
$g(t)$ is nonincreasing in $t$ gives
\begin{equation}
\label{harnackinequality}
R(x_2,t_2)\geq
\exp\left(-\frac{d_{t_1}^2(x_1,x_2)}{2(t_2-t_1)}\right)R(x_1,t_1).
\end{equation}
whenever $t_1 \: < \: t_2$ and $x_1, x_2 \in M$.
(If $n = 2$ then one can replace $\frac{d_{t_1}^2(x_1,x_2)}{2(t_2-t_1)}$ by
$\frac{d_{t_1}^2(x_1,x_2)}{4(t_2-t_1)}$.) In particular, if
$R(x_2, t_2) \: = \: 0$ for some $(x_2, t_2)$ then $g(t)$ must be flat
for all $t$.



\section{Alexandrov spaces} \label{Alexandrov}

We recall some facts about Alexandrov spaces
(see \cite[Chapter 10]{Burago-Burago-Ivanov},
\cite{Burago-Gromov-Perelman}).
Given
points $p,x,y$ in a nonnegatively curved Alexandrov space,
we let $\cangle_p(x,y)$ denote the comparison angle at $p$,
i.e. the angle of the Euclidean
comparison triangle at the vertex corresponding to $p$.

The Toponogov splitting theorem says that if $X$ is a
proper nonnegatively curved Alexandrov space which contains a line, then $X$
splits isometrically as a product $X=\R\times Y$, where
$Y$ is a proper, nonnegatively curved Alexandrov space
\cite[Theorem 10.5.1]{Burago-Burago-Ivanov}.

Let $M$ be an $n$-dimensional Alexandrov space with nonnegative curvature,
$p\in M$,
and $\la_k\ra 0$. Then the sequence $(\la_kM,p)$ of pointed spaces
Gromov-Hausdorff converges to the Tits cone $\ctits M$
(the Euclidean cone over the Tits boundary $\tits M$)
which is a nonnegatively curved Alexandrov
space of dimension $\leq n$
\cite[p. 58-59]{Ballmann-Gromov-Schroeder}.
If the Tits cone splits isometrically
as a product $\R^k\times Y$, then $M$ itself splits off
a factor of $\R^k$; using triangle comparison, one finds $k$
orthogonal lines passing through a basepoint, and applies the
Toponogov splitting theorem.

Now suppose that $x_k\in M$ is a sequence with $d(x_k,p)\ra\infty$
and $r_k \in \R^+$  is a sequence with
$\frac{r_k}{d(x_k,p)}\ra 0$.  Then the sequence
$(\frac{1}{r_k} M,x_k)$
subconverges to a pointed Alexandrov
space $(N_\infty,x_\infty)$ which splits off a line.  To see this,
observe that since
$(\frac{1}{d(x_k,p)}M,p)$
converges to a cone,
we can find a
sequence $y_k\in M$ such that $\frac{d(y_k,x_k)}{d(x_k,p)}\ra 1$,
and $\cangle_{x_k}(p,y_k)\ra\pi$.
Observe that for any $\rho<\infty$,
we can find sequences $p_k\in \ol{px_k}$, $z_k\in\ol{x_ky_k}$
such that $\frac{d(x_k,p_k)}{r_k}\ra \rho$,
$\frac{d(x_k,z_k)}{r_k}\ra \rho$, and by monotonicity of comparison
angles
\cite[Chapter 4.3]{Burago-Burago-Ivanov}
we will have $\cangle_{x_k}(p_k,z_k)\ra \pi$.  Passing
to the Gromov-Hausdorff limit, we find $p_\infty, z_\infty\in N_\infty$
such that $d(p_\infty,x_\infty)=d(z_\infty,x_\infty)=\rho$ and
$\cangle_{x_\infty}(p_\infty,z_\infty)=\pi$.  Since this
construction applies for all $\rho$, it follows that $N_\infty$
contains a line passing through $x_\infty$.
Hence, by the Toponogov splitting theorem, it is isometric
to a metric product $\R\times N'$ for some Alexandrov space
$N'$.

If $M$ is a complete Riemannian manifold of nonnegative sectional
curvature and $C\subset M$ is a
compact connected domain with weakly convex boundary then
the subsets $C_t= \{x\in C\mid d(x,\partial C)\geq t\}$
are convex in $C$
\cite[Chapter 8]{Cheeger-Ebin}.
If the second fundamental form of
$\partial C$ is $\geq \frac{1}{r}$ at each point
of $\partial C$, then for all $x\in C$ we have $d(x,\partial C)\leq r$,
since the first focal point of $\partial C$ along any inward pointing
normal geodesic occurs at distance $\leq r$.

Finally, we recall the statement of the Bishop-Gromov volume
comparison inequality.  Suppose that $M$ is an
$n$-dimensional
Riemannian manifold.
Given $p \in M$ and $r_2 \ge r_1 > 0$,
suppose that $B(p, r_2)$ has compact closure in $M$ and that the
sectional curvatures of $B(p, r_2)$ are bounded below by
$K \in \R$. Then
\begin{equation} \label{BG}
\frac{\vol(B(p, r_2))}{\vol(B(p, r_1))} \: \le \:
\begin{cases}
\frac{\int_0^{r_2} \sin^{n-1}(kr) \: dr}{\int_0^{r_1} \sin^{n-1}(kr) \: dr}
& \text{ if } K = k^2, \\ \\
\frac{r_2^n}{r_1^n}
& \text{ if } K = 0, \\ \\
\frac{\int_0^{r_2} \sinh^{n-1}(kr) \: dr}{\int_0^{r_1} \sinh^{n-1}(kr) \: dr}
& \text{ if } K =  -  k^2.
\end{cases}
\end{equation}
(If $K = k^2$ with $k>0$ then we restrict to $r_2 \le \frac{\pi}{k}$.)
The same inequality holds if we just assume that
$B(p, r_2)$ has Ricci curvature bounded below by
$(n-1) \: K$. Equation (\ref{BG}) also holds if
$M$ is an Alexandrov space and $B(p, r_2)$ has Alexandrov curvature
bounded below by $K$.

\section{Finding balls with controlled curvature} \label{pointpicking}

\begin{lemma}[existence for curvature]
  \label{lem:kl-finding-balls-controlled-curvature-existence-curvature}
  \dcref{Lemma H.1}
  \group{Kleiner--Lott Geometric Appendices}
  \source{kleiner-lott}

\label{selection}
Let $X$ be a Riemannian manifold
with $R \ge 0$
and suppose $\ol{B(x,5r)}$
is a compact subset of $X$.  Then there is a ball
$B(y,\bar r)\subset B(x,5r)$, $\overline{r} \le r$,
such that $R(z)\leq 2R(y)$
for all $z\in B(y,\bar r)$ and $R(y)\bar r^2\geq R(x)r^2$.
\end{lemma}
\begin{proof}
Define sequences $x_i\in B(x,5r)$, $r_i>0$
inductively as follows.  Let $x_1=x$, $r_1=r$.
For $i>1$, let $x_{i+1}=x_i$, $r_{i+1}=r_i$ if
$R(z)\leq 2R(x_i)$ for all $z\in B(x_i,r_i)$;
otherwise let $r_{i+1}=\frac{r_i}{\sqrt{2}}$, and let
$x_{i+1}\in B(x_i,r_i)$ be a point such that
$R(x_{i+1})>2R(x_i)$.  The sequence of balls
$B(x_i,r_i)$ is contained in $B(x,5r)$, so the
sequences $x_i,\,r_i$ are eventually constant,
and we can take $y=x_i$, $\bar r=r_i$ for large
$i$.
\end{proof}

There is an evident spacetime version of the lemma.

\section{Statement of the geometrization conjecture} \label{geom}

Let $M$ be a connected orientable closed (= compact boundaryless) $3$-manifold.
One formulation of the geometrization conjecture says that
$M$ is the connected sum of closed $3$-manifolds $\{M_i\}_{i=1}^n$, each of
which admits a codimension-$0$ compact submanifold-with-boundary $G_i$
so that
\begin{itemize}
\item $G_i$ is a graph manifold
\item $M_i - G_i$ is hyperbolic, i.e. admits a
complete Riemannian metric with constant negative sectional curvature
and finite volume
\item Each component $T$ of $\partial G_i$ is an incompressible torus in $M_i$,
i.e. with respect to a basepoint $t \in T$, the induced map
$\pi_1(T, t) \rightarrow \pi_1(M_i, t)$ is injective.
\end{itemize}
We allow $G_i = \emptyset$ or $G_i = M_i$.

A reference for graph manifolds is
\cite[Chapter 2.4]{Matveev}. The definition is as follows.
One takes a collection $\{P_i\}_{i=1}^N$ of pairs of pants
(i.e. closed $2$-disks with two balls removed)
and a collection of closed $2$-disks $\{D^2_j\}_{j=1}^{N^\prime}$. The
$3$-manifolds $\{S^1 \times P_i\}_{i=1}^N
\cup \{S^1 \times D^2_j\}_{j=1}^{N^\prime}$ have toral boundary components.
One takes an even number of these tori, matches them in pairs by
homeomorphisms, and glues $\{S^1 \times P_i\}_{i=1}^N
\cup \{S^1 \times D^2_j\}_{j=1}^{N^\prime}$ by these homeomorphisms.
The
resulting $3$-manifold $G$ is a graph manifold, and all
graph manifolds
arise in this way.
We will assume that the gluing homeomorphisms are such that
$G$ is orientable.
Clearly the boundary of $G$, if nonempty, is a
disjoint union of tori. It is also clear that the result of
gluing two graph manifolds along some collection of boundary tori is
a graph manifold. The connected sum of two $3$-manifolds is a
graph manifold if and only if each factor is a graph manifold
\cite[Proposition 2.4.3]{Matveev}.

The reason to require incompressibility of the tori $T$ in the
statement of the geometrization conjecture is to
exclude phony decompositions, such as writing $S^3$ as the
union of a solid torus and a hyperbolic knot complement.

A more standard version of the geometrization conjecture uses
some facts from $3$-manifold theory
\cite{Scott}. First $M$ has a Kneser-Milnor decomposition as
a connected sum of uniquely defined prime factors. Each prime factor is
$S^1 \times S^2$ or is irreducible, i.e. any embedded $S^2$ bounds
a $3$-ball. If $M$ is irreducible then it has a JSJ decomposition, i.e.
there is a minimal collection
of disjoint incompressible
embedded tori $\{T_k\}_{k=1}^K$ in $M$, unique up
to isotopy, with the property that
if $M^\prime$ is the metric completion of a component of
$M - \bigcup_{k=1}^K T_k$ (with respect to an induced
Riemannian metric from $M$) then
\begin{itemize}
\item $M^\prime$ is a Seifert $3$-manifold or
\item $M^\prime$ is non-Seifert and any embedded incompressible torus in
$M^\prime$ can be isotoped into $\partial M^\prime$.
\end{itemize}
The second version of the geometrization conjecture reduces to
the conjecture that in the latter case, the interior
of $M^\prime$ is hyperbolic. Thurston proved that this is true when
$\partial M^\prime \neq \emptyset$.
The reason for the word ``geometrization''
is explained in \cite{Scott,Thurston}.

An orientable Seifert $3$-manifold is a graph manifold
\cite[Proposition 2.4.2]{Matveev}. It follows that
the second version of the geometrization conjecture implies the
first version. One can show directly that any graph manifold
is a connected sum of prime graph manifolds, each of which can be split
along incompressible tori to
obtain a union of Seifert manifolds \cite[Proposition 2.4.7]{Matveev},
thereby showing the equivalence of the two versions.
